\section{Monday 09/30/2024}
\subsection{Review}
\subsubsection{Convex functions and sets}
\begin{equation}
  \frac{1}{2}x^T Q x + c^T x
\end{equation}
is the quadratic form and this function is convex as long as $Q$ is positive semidefinite.

\framebox[\linewidth]{
    \begin{minipage}{\dimexpr\linewidth-2\fboxsep-2\fboxrule\relax}
        A function $h: \mathbb{R}^n \to \mathbb{R}$ is said to be convex if and only if it satisfies 
        \begin{equation}
          \lambda h(x_1) + (1-\lambda)h(x_2) \geq h(\lambda x_1 + (1-\lambda)x_2)
        \end{equation}  
        for any $\lambda \in [0,1]$ and any $x_1, x_2 \in \textbf{dom } h$ 
    \end{minipage}
}
\\ \\ 
\framebox[\linewidth]{
    \begin{minipage}{\dimexpr\linewidth-2\fboxsep-2\fboxrule\relax}
        A set $\Omega \subseteq \mathbb{R}^n $ is said to be convex if and only if it satisfies 
        \begin{equation}
          \lambda x_1 + (1-\lambda)x_2 \in \Omega
        \end{equation}  
        for any $\lambda \in [0,1]$ and any $x_1, x_2 \in \Omega$ 
    \end{minipage}
}


\subsection{Convexity}
\subsubsection{Properties of a convex function}
For a twice-differential function to be convex, it's Hessian matrix must be positive semidefinite. One definition for positive semidefinite is that all eigenvalues must be non-negative, meaning that the stretch of the matrix is positive or none in the direction of each eigenvector. Below we define what the hessian of a function $f$ looks like.
\begin{align}
  \nabla^2 f(x) = 
  \begin{bmatrix}
     \frac{d^2 f(x)}{d x_1^2 } \frac{d^2 f(x)}{d x_1 x_2 } \dots \frac{d^2 f(x)}{d x_1 x_n} \\
     \frac{d^2 f(x)}{d x_2 x_1 } \frac{d^2 f(x)}{d x_2^2} \dots \frac{d^2 f(x)}{d x_2 x_n} \\
     \vdots\\
     \frac{d^2 f(x)}{d x_n x_1 } \frac{d^2 f(x)}{d x_n x_2} \dots \frac{d^2 f(x)}{d x_n^2}    
  \end{bmatrix}
\end{align}
\\ \\
\framebox[\linewidth]{
    \begin{minipage}{\dimexpr\linewidth-2\fboxsep-2\fboxrule\relax}
        \textbf{Note:} Diagonal matrices' eigenvalues are the values along the diagonal.
    \end{minipage}
}

\subsubsection{Examples}
\begin{itemize}
  \item The 1D function $f(x) = \frac{1}{2}x^2$ is a convex function because its hessian is $\nabla^2 f(x) = 1$. Since this is a positive value and this matrix is positive semidefinite, the function $f$ is convex. 
  \item The 1D function $f(x) = x^3$ has a hessian of $\nabla^2 f(x) = 6x$. This function is convex on $\mathbb{R}_{++}$ and concave otherwise (both at x =0).
  \item The 1D function $f(x) = \frac{1}{x}$ is convex on the domain $\mathbb{R_{++}}$. The hessian of this function is $\nabla^2 f(x) = \frac{2}{x^3}$, which is positive on $\mathbb{R_{++}}$
  \item The 1D function $f(x) = -log(x)$ is convex on the domain $\mathbb{R_{++}}$. The hessian of this function is $\nabla^2 f(x) = \frac{1}{x^2}$
\end{itemize}

\subsubsection{Linear regression}
Attempting to verify if the least squares equation is convex.
\begin{equation}
  \min_x \frac{1}{2n} \sum_{i=1}^n (a_i^T x - b_i)^2
\end{equation}
\begin{gather}
  \nabla f(x) = \frac{1}{2n} \sum_{i=1}^n \frac{d }{d x} (a_i^T x -b_i)^2 \\
  = \frac{1}{2n} \sum_{i=1}^n (a_i^T x -b_i)a_i \\
  \nabla^2 f(x) = \frac{1}{2n} \sum_{i=1}^n \frac{d }{d x} a_i^T x a_i - b_i a_i \\
  = \frac{1}{2n} \sum_{i=1}^n a_i^T a_i
\end{gather}
This is convex since $a_i^T a_i$ is a symmetric positive semi-definite matrix.

\section{Wednesday 10/02/2024}
\subsection{Convex Functions}
\subsubsection{Verifying convexity}
\begin{itemize}
  \item $\frac{1}{2n} \sum_{i=1}^n (a_i^T x - b_i)^2$
  \item $a^T x + b$
  \item $-\log(x)$
  \item $\frac{1}{x}, x>0$
  \item $\frac{1}{x}, x \neq 0$
  \item $\frac{1}{x}, x<0$
  \item $g(x) + h(x)$
  \item $\frac{g(x)}{h(x)}$
  \item $g(h(x))$
  \item $\max\{f_1(x),f_2(x)\}$
\end{itemize}
\subsubsection{composition rule for convexity}
A logical conclusion for a composition rule $f(x) = g(h(x))$ would be that $f$ is convex if both $g$ and $h$ are convex. However, this does not work because we can use the example $g(x) = -x$ and $h(x) = -\log(x)$. The composite would be $f(x) = \log(x)$ which is concave, not convex. Therefore in order to construct a set of rules for establishing convexity, we can take the second derivate of a composite
\begin{gather}
  f(x) = g(h(x)) \\ 
  \nabla f(x) = h^\prime(x) g^\prime(h(x)) \\
  \nabla^2 f(x) = h^{\prime \prime}(x) g^\prime(h(x)) + h^\prime(x)^2 g^{\prime \prime}(h(x))
\end{gather}
$g$ needs to be convex, and either $h$ is convex and $g$ is non-decreasing or $h$ is concave and $g$ is non-increasing.

\subsection{Convex Sets}
If a function $f$ is convex, then the set
\begin{equation}
  \Omega = \{ x | f(x) \leq c \}
\end{equation}
is convex. The notation above means "The set of all $x$ given that a function $f(x)$ is less than $c$"
\begin{gather}
  \lambda f(x_1) + (1-\lambda) f(x_2) \geq f(\lambda x_1 + (1-\lambda)x_2) \\
  \text{since } f(x_1), f(x_2) \leq c \\
  \lambda c + (1-\lambda) c \geq f(\lambda x_1 + (1-\lambda)x_2) \\ 
  c \geq f(\lambda x_1 + (1-\lambda)x_2)
\end{gather}
The last line implies that any point on the line $f$ that is between two arbitrary points $x_1,x_2 \in \Omega$ is less than $c$. \\
What about the set 
\begin{equation}
  \hat{\Omega} = \{x | f(x) = c\} 
\end{equation}

This set is not convex. Equality constraints are convex if $f$ is a linear function.

\section{Friday 10/04/2024}
\subsection{Preserving Convexity}
\subsubsection{Non-negative weighted sum}
if $f_1(x), f_2(x)$ are convex functions. Then their non-negative weighted sum 
\begin{equation}
  g(x) = \alpha f_1(x) + \beta f_2(x)
\end{equation}
is convex
\subsubsection{affine composition}
If $f(x)$ is convex and $x$ can be written as the affine composition $x=Ay +b$, then $f(Ay+b)$ is also convex

\subsection{Applications}
Many real-world functions are not convex, so attempting to build convex representations can lose important signals and aspects of the original function. In the application of image reconstruction, taking an image that has been blurred and distorted and reconstructing it back into the original resolution of the image.
\subsubsection{Application 1: Upscaling images}
One application of convex optimization is through upscaling images that are noisy and blurry in attempts to recover the original detail of the image.
\begin{figure}[htbp]
  \centering
  \begin{subfigure}[b]{0.49\textwidth}
    \centering
    \includegraphics[width=\textwidth]{images/background_pixellated.png}
    \caption{Blurry and noisy image}
    \label{fig:background_pixellated}
  \end{subfigure}
  \hfill
  \begin{subfigure}[b]{0.49\textwidth}
    \centering
    \includegraphics[width=\textwidth]{images/background_original.png}
    \caption{Original image being recovered}
    \label{fig:background_original}
  \end{subfigure}
  \caption{Comparison between the blurry and noisy image (left) and the original image (right).}
  \label{fig:comparison_images}
\end{figure}


\subsubsection{Radiotherapy treatment planning}
Convex optimization can be useful in radiotherapy treatment. Radiation damages the DNA of tumor cells. Properly optimized radiotherapy machines aim to maximize the amount of tumor damaged by the machine while minimizing the damage to other cells. The stochasticity in the partical/xray process introduces non-linearity into the problem. Some constraints are in the problem because there are critical organs that must be avoided in the process. 
\begin{equation}
  \begin{aligned}
    \min f(D) \\
    s.t. D = F(x) \\
    x \geq 0 \\
    \text{other dose constraints}
  \end{aligned}
\end{equation}
Where $F$ is a function that maps the source of radiation to the distribution of dosage. $f$ is better dose conforming given dose prescription. $x$ is a vector with intensity, angle, and shape attributes. $x$ is the intensity of radiation from its sources. Each beam is separated into a grid, and then collected into the variable $x$ where each entry in the vector corresponds to an intensity of a beam at a certain beamlet. $D$ is dose distribution. The amount of energy on each region (voxel, a small cube unit) of the human body. 
\begin{align}
  D = 
  \begin{bmatrix}
     1.4 \\
     3.3 \\
     5.2
  \end{bmatrix}
\end{align} 
Here, that means the region is partitioned into three voxels, and the amount of dose in each region is 1.4, 3.3, and 5.2 respectively.

\begin{align}
  x = 
  \begin{bmatrix}
     0 \\
     3 \\
     5 \\
     0 \\
     3 \\
     2 \\
  \end{bmatrix}
\end{align} 
Here, the x is a flattened matrix that contains the intensity and direction of a bunch of beamlets collapsed into one vector.
\\ \\
Problem formulation:
\begin{itemize}
  \item target: 95\% of volume $\geq$ 72
  \item critical organ 1 average $\leq$ 10
  \item Critical organ 2 max dose $\leq$ 72
\end{itemize}